% Transcription — fonds Alexandre Grothendieck, shelfmark 132, batch 2 (pages 21-31).
% Established from the facsimile archives/batches/132/batch-02.pdf, cut from a
% copy of 132.pdf fetched through the project's own relay
% (https://grothendieck-relay.onrender.com/source/132.pdf), not uploaded by the
% user, per the skill transcribe-grothendieck. The batch file carries no cover
% sheet: PDF page k is page 20+k of the folder, and the archivists' pencil
% numbers bottom-left agree on every leaf (« 21 » ... « 31 »). The folder is
% thirty-one pages; this is the second and last batch (pages 1-20 are batch 1,
% transcribed separately and in parallel, so the notation below was fixed from
% these pages alone).
%
% Pass: Opus 5.5 (claude-opus-5-5), 2026-09-26 — first pass, unchecked against the pages by a human.
%
% Dating. The inventory dates the folder « 1974 » (copied verbatim into
% \dating), which is the date of the offprints. The folder sits in the
% inventory group « Descentes » (126-128, which also takes 129-133), dated
% « [avant 1970] »: the group's date and the folder's own do not agree.
% Recorded as an observation only. The manuscript leaves carry no date.
%
% Pages transcribed from his hand: 21, 22, 23 (blue ink, fast drafting; the
% formulas and diagrams carry, part of the connective prose does not).
% Page skipped: 24, a blank cream sheet.
% Pages summarised: 25-31, the « tirés à part (1974) » of the inventory — two
% printed Notes aux Comptes rendus by Yves Ladegaillerie, « Classification
% topologique des plongements des 1-complexes compacts dans les surfaces »:
% pp. 25-28 are the second Note (C. R. Acad. Sc. Paris, t. 279, 22 juillet
% 1974, Série A, pp. 129-132, §§ 7-13), pp. 29-31 the first (t. 278, 27 mai
% 1974, pp. 1401-1403, §§ 1-5), filed in that order. They are another
% author's printed text and are not re-set: following folder 56, each page
% gets one short French \note saying which paragraphs it carries. No mark of
% his hand relating to the text was found on them. The only manuscript trace
% is a pencil sketch at the foot of page 28 (the blank lower half of the last
% page of the second Note), which looks like a pan or bowl and has no
% connection with the mathematics; it cannot be attributed and is not
% transcribed. The author's name is a published byline and is given as such.
%
% His pagination: none on these leaves. Page 22 opens a paragraph marked
% « (C) » in a circle, written over a struck « (D) »; page 20 (batch 1) ends
% on a paragraph « (B') », so these letters are his own sectioning, and the
% conditions (A), α), β), γ) and the data a), b), c) referred to on pages
% 22-23 are set up in batch 1.
%
% What the batch is about. The manuscript leaves have nothing to do with the
% offprints: they continue an argument on extensions of topological groups
% 1 -> N -> G -> H -> 1, with H~ the universal covering of H, pi = pi_1(H)
% (1 -> pi -> H~ -> H -> 1), gothic G a discrete quotient of H~ by H~_0,
% and a homomorphism phi : pi -> z(N) into the centre of N.
%   21     end of (B'): pi_1 G = Ker(pi_1(H) -> pi_0 N), the diagram through
%          pi_0(pi) = pi (« car pi discret », by the homotopy exact sequence),
%          the exact sequence pi_1 H -> N -> pi_0 G -> pi_0 H -> 1, so
%          pi_0 G an extension of pi_0 H by Coker(pi_1 H -> N); then, via
%          the extension G_0 of gothic G by N and H -> gothic G, an
%          extension G_1 of pi_0 H by N and pi_0 G = Coker(pi_1 H -> G_1).
%          A margin word, read « faux » with doubt, brackets that last claim.
%   22-23  (C): functoriality in a homomorphism u : H' -> H. Pulled-back
%          sequences (b'), (a'), and G' = G x_H H'; the pullback to H~'_0 is
%          strictly trivialised, so G' is determined by 1°) an extension of
%          gothic G' by N and 2°) a homomorphism phi' : pi -> G~'_0
%          centralising N and lifting pi -> gothic G'; G~'_0 is the pullback
%          of G~_0, phi' is induced by phi. Conditions beta), gamma) pass to
%          the pulled-back data but alpha) (H~'_0 connected and simply
%          connected) does not; the closing sentence, on pi_0 G~', is
%          cancelled by long diagonal strokes.
%
% Notation fixed for the batch: \mathfrak{G}, \mathfrak{G}' for his gothic G
% (a curled capital, closed loop at the left); \mathfrak{z}(N) for the curly
% « 3 » he writes for the centre of N (hand.md § 4); \widetilde{H},
% \widetilde{H}_0, \widetilde{H}'_0, \widetilde{G}_0, \widetilde{G}'_0 as he
% writes them (the prime often sits over the index 0); \varphi drawn across
% the shaft of its arrow is set as an arrow label. His dotted arrows are
% `dashed`; the wavy arrow of page 22 is set plain, with a note. Square
% brackets in the text are his. Counts: \ill 13, \uncertain 12.
\documentclass[11pt,a4paper]{article}
% Le préambule commun à toutes les transcriptions du fonds Grothendieck.
%
% Il tient en une idée : **l'incertitude se compose, elle ne se résout pas.**
% Une transcription qui lisse un mot douteux a détruit la seule chose pour
% laquelle on la faisait — un lecteur doit pouvoir savoir, à chaque endroit,
% ce qui était sur le feuillet et ce qui a été ajouté. Les six macros qui
% suivent sont donc l'appareil critique, et non de la décoration.
%
% Les mêmes commandes sont reconnues par `scripts/render.mjs`, qui produit la
% vue de lecture du site. Une macro ajoutée ici doit l'être aussi là-bas,
% faute de quoi le rendu échouera bruyamment — ce qui est voulu : un rendu
% silencieusement faux est pire qu'un rendu absent.

\ProvidesPackage{grothendieck}[2026/08/08 Transcriptions du fonds Grothendieck]

\usepackage{amsmath,amssymb,amsthm}
\usepackage{mathrsfs}
\usepackage{stmaryrd}
% Les diagrammes commutatifs sont partout dans ce fonds : c'est la langue
% même de la géométrie algébrique de Grothendieck, et une transcription qui
% les rendrait en prose ne transcrirait rien.
\usepackage{tikz-cd}
% Français par défaut — c'est la langue des feuillets — mais l'anglais est
% chargé aussi : la traduction et le résumé sont des documents anglais, et la
% typographie française (espaces avant les deux-points) y serait une faute.
% Ces documents basculent par \selectlanguage{english} en tête de corps.
\usepackage[english,french]{babel}
\usepackage{fontspec}
\usepackage{geometry}
\geometry{a4paper,margin=25mm,marginparwidth=28mm}
\usepackage{marginnote}
\usepackage{xcolor}
% ulem, not soul. `soul`'s \st and \ul are built on a character-by-character
% scan that XeLaTeX's font machinery breaks: the document fails with an
% undefined control sequence rather than degrading. ulem does the same two jobs
% and is Unicode-engine safe. `normalem` stops it hijacking \emph, which the
% transcriptions use for Grothendieck's own emphasis.
\usepackage[normalem]{ulem}
\usepackage{hyperref}
\hypersetup{colorlinks=true,linkcolor=black,urlcolor=black,pdfborder={0 0 0}}

% --- Métadonnées du lot -----------------------------------------------------
% Reprises telles quelles par le script de rendu, qui les affiche en tête de
% la vue de lecture. Elles fixent ce que le fichier prétend couvrir : sans
% elles, une transcription flotte sans point d'attache dans un fonds de seize
% mille feuillets.

\newcommand{\folder}[1]{\gdef\grothFolder{#1}}
\newcommand{\batch}[1]{\gdef\grothBatch{#1}}
\newcommand{\leaves}[2]{\gdef\grothFirst{#1}\gdef\grothLast{#2}}
\newcommand{\foldertitle}[1]{\gdef\grothFoldertitle{#1}}
\gdef\grothFolder{?}\gdef\grothBatch{?}\gdef\grothFirst{?}\gdef\grothLast{?}\gdef\grothFoldertitle{}

% --- Le résumé --------------------------------------------------------------
%
% Il ouvre la lecture modernisée, sous le titre « Résumé », et il est composé
% en petit. Ce n'est pas un ornement : c'est le seul passage du document qui
% ne soit pas des mathématiques, et un lecteur doit pouvoir le reconnaître
% avant de l'avoir lu — puis le sauter s'il connaît déjà le sujet, ou s'y
% arrêter s'il ne le connaît pas. Le filet qui le suit marque la reprise du
% corps.

\newenvironment{resume}
  {\small\setlength{\parskip}{0.45em}}
  {\par\medskip\hrule\medskip}

% --- Mots-clés ---------------------------------------------------------------
%
% En anglais, en clôture du résumé de la lecture modernisée : le vocabulaire
% moderne sous lequel chercher ce que les feuillets construisent. Le manifeste
% du site (`npm run manifest`) les extrait pour étiqueter la cote entière —
% cette ligne est la source unique des tags, il n'y a pas de fichier à côté.
\newcommand{\keywords}[1]{\par\smallskip\noindent
  {\footnotesize\textsc{Keywords} --- \textit{#1}}\par}

% --- L'appareil critique ----------------------------------------------------

% Le numéro de feuillet, en marge. C'est l'ancre qui permet de régler un
% désaccord sur une formule en regardant *un* feuillet plutôt qu'une chemise
% de six cents. Le site s'en sert aussi pour tourner les pages du fac-similé
% quand on fait défiler la transcription.
\newcommand{\leaf}[1]{%
  \marginnote{\footnotesize\color{gray!70}\textsf{#1}}%
  \label{leaf:#1}\ignorespaces}

% Illisible : on ne devine pas. Un mot inventé qui se lit comme les autres est
% le pire résultat possible.
\newcommand{\ill}{\textcolor{red!60!black}{[\,\ldots\,]}}

% Lecture douteuse : le mot est donné, et signalé comme incertain.
\newcommand{\uncertain}[1]{\uline{#1}}

% Ajout de l'éditeur — une lettre manquante, un mot sous-entendu.
\newcommand{\add}[1]{\textcolor{blue!50!black}{[#1]}}

% Biffé par Grothendieck lui-même. Ce qu'il a rayé fait partie du document :
% une rature dit souvent où la pensée a changé de direction.
\newcommand{\struck}[1]{\sout{#1}}

% Note du transcripteur, sur l'état matériel du feuillet.
\newcommand{\note}[1]{\par\smallskip{\small\color{gray!60!black}\textit{[#1]}}\par\smallskip}

% Note marginale de Grothendieck, distincte du corps du texte.
\newcommand{\marginal}[1]{\par\smallskip\noindent\hspace{1em}%
  {\small\color{gray!60!black}#1}\par\smallskip}

% --- Titre ------------------------------------------------------------------

\AtBeginDocument{%
  \begin{center}
    {\large\bfseries Fonds Alexandre Grothendieck --- cote n\textsuperscript{o}~\grothFolder}\\[2pt]
    {\itshape\grothFoldertitle}\\[4pt]
    {\small feuillets \grothFirst--\grothLast\ (lot \grothBatch)}\\[2pt]
    {\footnotesize\color{gray!60!black}Transcription établie d'après le fac-similé numérisé
     par l'Université de Montpellier.\\
     Les lectures incertaines sont soulignées, les illisibles notées [\,\ldots\,],
     les ajouts entre crochets.}
  \end{center}
  \vspace{1em}}

\endinput


\folder{132}
\batch{2}
\pages{21}{31}
\dating{1974}
\watermark{Édition de démonstration}
\foldertitle{Plongements dans surfaces (Ladegaillerie) : notes manuscrites (s.d.), tirés à part (1974).}

\begin{document}

\page{21}
\note{suite du paragraphe (B') de la page 20}

\[
  \pi_1 G = \operatorname{Ker}\bigl(\pi_1(H) \longrightarrow \pi_0 N\bigr)
\]
\note{au-dessus de cette flèche, « (connu) », avec un petit trait qui la
désigne ; la formule est complétée en diagramme :}

\begin{tikzcd}
  \pi_1(H) \arrow[r] \arrow[dd] & \pi_0 N \arrow[d, no head, "\parallel" description] \\
  & N \\
  \pi_0(\pi) = \pi \arrow[r, "\varphi"'] & \mathfrak{z}(N) \arrow[u, hook]
\end{tikzcd}

\note{un double trait oblique part de $\pi$ vers l'inclusion
$\mathfrak{z}(N) \subset N$ ; une flèche courbe renvoie de « car $\pi$
discret » à l'égalité $\pi_0(\pi) = \pi$}

\marginal{car $\pi$ discret}

[\uncertain{Démonstr.} : suite exacte d'homotopie pour ext.
$1 \to \pi \to \widetilde{H} \to H \to 1$]
\note{entre crochets, à gauche, relié par un trait à la flèche verticale
$\pi_1(H) \to \pi_0(\pi)$}

\struck{\ill{}} et on a

\begin{tikzcd}
  \pi_1 H \arrow[r] \arrow[d, "\partial"'] & N \arrow[r] & \pi_0 G \arrow[r] & \pi_0 H \arrow[r] & 1 \\
  \pi \arrow[r, "\varphi"] & \mathfrak{z}(N) \arrow[u, hook] & & &
\end{tikzcd}

\note{au-dessus de la flèche $\pi_1 H \to N$, un mot désigné par un petit
trait, \uncertain{connu}}

i.e. $\pi_0 G$ \struck{\ill{}} est une extension de $\pi_0 H$ par
$\operatorname{Coker}(\pi_1 H \to N)$]. Or on a une extension $G_0$ de
$\mathfrak{G}$ par $N$ (définie car $\pi \subset \widetilde{H}_0$)
\uncertain{d'où}, grâce à $H \to \mathfrak{G}$ (i.e.
$\pi_0 H \to \mathfrak{G}$ ($= \pi_0 \mathfrak{G}$) car $\mathfrak{G}$
discret) une extension $G_1$ de $\pi_0 H$ par \struck{\ill{}} $N$ ;
\note{« (définie car $\pi \subset \widetilde{H}_0$) » est ajouté sous la
ligne ; le crochet fermant après $N$ n'a pas d'ouvrant sur la page}

\marginal{\uncertain{faux}}
\note{mot en marge gauche, suivi d'une croix, avec un grand trait vertical
qui embrasse la phrase suivante et sa formule}

on vérifie alors que
\[
  \pi_0 G \simeq \operatorname{Coker}\bigl(\pi_1 H \longrightarrow G_1\bigr)
\]

\begin{tikzcd}
  \pi_1 H \arrow[r] \arrow[dd, "\partial"'] & G_1 \\
  & N \arrow[u, hook] \\
  \pi \arrow[r, "\varphi"] & \mathfrak{z}(N) \arrow[u, hook]
\end{tikzcd}

\page{22}
\struck{(D)} (C) Soit, sous les conditions de (A),
\[
  H' \xrightarrow{\;u\;} H
\]
un hom. de groupes topologiques. On en \uncertain{déduit} par
\uncertain{images inverses} une extension

\begin{tikzcd}
  (b') & 1 \arrow[r] & \pi \arrow[r] \arrow[d, no head, "\parallel" description] & \widetilde{H}' \arrow[r] \arrow[d] & H' \arrow[r] \arrow[d, "u"] & 1 \\
  (b) & 1 \arrow[r] & \pi \arrow[r] & \widetilde{H} \arrow[r] & H \arrow[r] & 1
\end{tikzcd}

et désignant par $\widetilde{H}'_0$ l'image inverse de $\widetilde{H}_0$ par
$\widetilde{H}' \to \widetilde{H}$, \uncertain{on trouve}, posant
$\mathfrak{G}' = \widetilde{H}' / \widetilde{H}'_0 \hookrightarrow \mathfrak{G}$,
une suite exacte

\begin{tikzcd}
  (a') & 1 \arrow[r] & \widetilde{H}'_0 \arrow[r] \arrow[d] & \widetilde{H}' \arrow[r] \arrow[d] & \mathfrak{G}' \arrow[r] \arrow[d, hook] & 1 \\
  (c) & 1 \arrow[r] & \widetilde{H}_0 \arrow[r] & \widetilde{H} \arrow[r] & \mathfrak{G} \arrow[r] & 1
\end{tikzcd}

\note{dans la ligne (a'), un symbole biffé avant $\widetilde{H}'_0$ ;
l'étiquette de la seconde ligne se lit « (c) », non « (a) » comme
l'appariement avec (a') le ferait attendre}

Considérons maintenant l'extension $G'$ de $H'$ par $N$ image inverse de $G$
par $H' \to H$

\begin{tikzcd}
  1 \arrow[r] & N \arrow[r] \arrow[d, no head, "\parallel" description] & G' \arrow[r] \arrow[d] & H' \arrow[r] \arrow[d] & 1 \\
  1 \arrow[r] & N \arrow[r] & G \arrow[r] & H \arrow[r] & 1
\end{tikzcd}

\note{les deux lignes portent à gauche une étiquette entre parenthèses (la première accentuée), dont la lettre, réduite à un petit trait, est illisible : \ill{}}

son image inverse par $\widetilde{H}'_0 \to H'$ est munie

\begin{tikzcd}
  \widetilde{H}'_0 \arrow[r] \arrow[d, dashed] & H' \arrow[d, dashed] \\
  \widetilde{H}_0 \arrow[r] & H
\end{tikzcd}

\note{les deux flèches verticales sont pointillées ; la flèche
$\widetilde{H}_0 \to H$ est ondulée}

de façon évidente d'une \uncertain{trivialisation stricte},
\struck{\ill{}} (grâce à celle de l'image inverse de $G$ par
$\widetilde{H}_0 \to H$) donc elle est déterminée par

1°) une ext. $\widetilde{G}'_0$ \struck{de} de $\mathfrak{G}'$ par $N$
\[
  1 \longrightarrow N \longrightarrow \widetilde{G}''_0 \longrightarrow
  \mathfrak{G}' \longrightarrow 1
\]
\note{dans la suite exacte, l'objet du milieu porte deux accents, là où la
ligne précédente n'en porte qu'un}

\page{23}
2°) un hom. $\pi \xrightarrow{\;\varphi'\;} \widetilde{G}'_0$ qui
\uncertain{centralise} $N$ et relevant $\pi \to \mathfrak{G}'$ [donc, si
$\pi \subset \widetilde{H}'_0$ d'où \struck{$\pi \to$}
$\pi \to \mathfrak{G}' \subset \mathfrak{G}$ trivial, un hom.
\[
  \pi \xrightarrow{\;\varphi'\;} \mathfrak{z}(N) \qquad ]
\]

Ceci posé, on vérifie immédiatement que l'extension $\widetilde{G}'_0$ est
l'image inverse de $\widetilde{G}_0$ par $\mathfrak{G}' \to \mathfrak{G}$,
\struck{et que} (donc $\widetilde{G}'_0 \subset \widetilde{G}_0$), et que
$\varphi'$ est induit par $\varphi$ [donc, si $\pi \subset \widetilde{H}_0$,
$\varphi = \varphi'$ en tant que hom. $\pi \to \mathfrak{z}(N)$]

Si les conditions \struck{\ill{}} $\beta$) \uncertain{$\gamma$)} sont vérifiées pour les
données \struck{\ill{}} b), c), il en est de même pour les données.
\struck{\ill{}} a'), b'), c') ($\pi$, $N$, $\mathfrak{G}'$ discret,
$\pi \subset \widetilde{H}'_0$) mais la condition $\alpha$ ($\widetilde{H}'_0$
connexe et simplement connexe \ill{} \uncertain{impliquée} par \ill{},
\ill{} que $\widetilde{H}'_0$ soit connexe~…) Mais comme
$\mathfrak{G}'$ est discret, \struck{\ill{}} il en est de \uncertain{même}
de $\mathfrak{G}'_0$, donc $G'_0$ \ill{} comme un quotient de
$\pi_0 \widetilde{G}'$
\note{les trois dernières lignes, depuis « Mais comme », sont barrées de
quatre longs traits obliques ; l'accent de $\widetilde{H}'_0$ dans la
condition $\alpha$ ressemble à un astérisque}

\section*{Les tirés à part}

\page{25}
\note{première page d'un tiré à part imprimé : Yves Ladegaillerie,
« Classification topologique des plongements des 1-complexes compacts dans
les surfaces », Note présentée par M. Henri Cartan, \emph{C. R. Acad. Sc.
Paris}, t. 279 (22 juillet 1974), Série A, p. 129 — la seconde des deux
Notes : §§ 7 (voisinages tubulaires, théorème 7.1, catégories TUB et TUBO),
8 (1-complexes circularisés, théorème 8.3) et début du § 9. Rien de sa main}

\page{26}
\note{p. 130 de la même Note : fin du § 9 (proposition 9.3), § 10
($\mathbf{Z}_2$-1-complexes circularisés, théorème et corollaire), § 11
(cocircularisations, théorème 11.4, terminologie des circuits et des
boucles). Rien de sa main}

\page{27}
\note{p. 131 de la même Note : fin du § 11, § 12 (classification des
plongements, proposition 12.4) et § 13 (plongements minimaux, genre de $K$,
théorème 13.5). Rien de sa main}

\page{28}
\note{p. 132 de la même Note : les trois références et l'adresse de l'auteur
(USTL, Montpellier) ; le reste de la page est vide. Rien de sa main qui
touche au texte}

\page{29}
\note{première page d'un second tiré à part : Yves Ladegaillerie, même titre,
Note présentée par M. Henri Cartan, \emph{C. R. Acad. Sc. Paris}, t. 278
(27 mai 1974), Série A, p. 1401 — la première des deux Notes : § 1
(catégorie des 1-complexes à isotopie près, théorème 1.4) et § 2
(préliminaires topologiques, 2.1-2.4). Rien de sa main}

\page{30}
\note{p. 1402 de la même Note : proposition 2.5, § 3 (voisinage tubulaire
standard d'un 1-complexe circularisé), § 4 (voisinage tubulaire d'un
plongement, théorème 4.1, corollaire 4.2) et début du § 5. Rien de sa main}

\page{31}
\note{p. 1403 de la même Note : fin du § 5, théorème de classification 5.2,
les quatre références et l'adresse de l'auteur. Rien de sa main}

\end{document}
